\documentclass[10pt, a5paper, twoside]{ltjsarticle}

%\usepackage{okumacro}
\usepackage{enumitem}
\usepackage{amssymb}
\usepackage{amsmath}
\usepackage{amsfonts}
\usepackage{amsthm}
\usepackage{mathrsfs}
\usepackage{bm}
\usepackage{xurl}
\usepackage{here}
\usepackage{graphicx}
\usepackage{wrapfig}
\usepackage{makeidx}
\usepackage{braket}
%\usepackage{ascmac}
\usepackage{fancyhdr}
\usepackage{tcolorbox}
\usepackage{tikz}
%\usepackage{otf}
\usepackage{luatexja-otf}
\usepackage{luatexja-ruby}
\usepackage[top=20truemm,bottom=20truemm,left=15truemm,right=15truemm]{geometry}
\usepackage{mathrsfs}

\pagestyle{fancy}
	\fancyhead{}
	\fancyhead[RE]{円城塔研究}
	\fancyhead[LO]{資料：コンメンタール「Your Heads Only」}
	\fancyhead[LE, RO]{\thepage}
	\fancyfoot{}
	\fancyfoot[LE, RO]{\footnotesize{The Journal of EnJoeToh Research, Vol.4, No.2, 2026} }

\tcbuselibrary{breakable, theorems, skins}

\newtcbtheorem{axm}{公理}{enhanced, breakable=true}{axm}
\newtcbtheorem[use counter from=axm]{dfn}{定義}{enhanced, breakable=true}{dfn}
\newtcbtheorem[use counter from=axm]{thm}{定理}{enhanced, breakable=true}{thm}
\newtcbtheorem[use counter from=axm]{sys}{系}{enhanced, breakable=true}{sys}
\newtcbtheorem[use counter from=axm]{exm}{例}{enhanced, breakable=true}{exm}
\newtcbtheorem[use counter from=axm]{hyp}{予想}{enhanced, breakable=true}{hyp}
\newtcbtheorem[use counter from=axm]{ths}{提唱}{enhanced, breakable=true}{ths}
\newtcbtheorem[use counter from=axm]{prn}{原理}{enhanced, breakable=true}{prn}
\newtcbtheorem[use counter from=axm]{clm}{主張}{enhanced, breakable=true}{clm}
\newtcbtheorem[use counter from=axm]{prb}{問題}{enhanced, breakable=true}{prb}
\newtcbtheorem[use counter from=axm]{prf}{証明}{enhanced, breakable=true, colback=white, attach boxed title to top left={xshift=2mm, yshift=-2mm}}{prf}

\newcommand{\hyphen}{\insertkanjiskip--\insertkanjiskip}

\newcommand{\cross}[1]{%
\mbox{\vbox{\kern 1pt\hbox{\vbox{\hrule
\kern 2pt\hbox{\ensuremath{\vphantom{b}#1}\kern 2pt}}\vrule\kern 1pt}}}\,}

\begin{document}

	~ %強制改行

	\begin{center}

		\Large{コンメンタール「Your Heads Only」}

		\vspace{3mm}

		\large{Commentary of short story \textit{Your Heads Only}}

		\vspace{3mm}

		\large{下村思游}

	\end{center}

	\vspace{3mm}

	\section{解説}

        頁数・行数は『Boy's Surface』紙版\cite{boys}に準拠した．

        \subsection{題名}

        元ネタは“読後焼却すべし”を意味する``Your Eyes Only''から．映画『007 ユア・アイズ・オンリー』からの連想も指摘される．headsは読者の頭と，チューリング機械の読み取り部（ヘッド）の両方を意味する．

        円城塔は，物語はテクストと読者（の脳神経構造上の電子の集合）との相互作用によって記述されると繰り返し主張しており，その代表例が本作である\footnote{「ムーンシャイン」『プロローグ』『文字渦』「for Smullyan」「ドルトンの印象法則」をはじめ，多くの作品に普遍的に見られる．}．テクストと読者から成る物語という系のモデルは，テクストをテープ，読者をヘッドとみなすことでチューリング機械として自然に解釈される．この発想は，生命系をチューリング機械の類似でモデルした池上機械からの連想であろう\footnote{池上機械は生命系そのものをテープ\hyphen ヘッド系としてモデルする一方，物語のモデルではテクストをテープ，生命系である読者をヘッドとしており，生命系がテープ\hyphen ヘッド系のどの部分に相当するかという点に相違がある．逆に，物語は生命系であると主張していると見ることも出来る．金子邦彦は，複雑系研究に最適な手法は物語であると主張している\cite{chaos}p.94．また，初期の複雑系の研究者で統計学者の伊庭幸人は，最新の科学理論は生命の理論になると主張する．}．ここで，チューリング機械は以下のように定義される\cite{sipser2}．

        \begin{dfn}{チューリング機械}{}

			\textbf{チューリング機械}は，7個組$(Q, \Sigma, \Gamma, \delta, q_0, q_{\mathrm{accept}}, q_{\mathrm{reject}})$である．ここで，$Q, \Sigma, \Gamma$はすべて有限集合であり，以下の性質をすべて満たす．

			\begin{enumerate}
				\item $Q$は状態の集合

				\item $\Sigma$は空白文字$␣$を含まない入力アルファベット

				\item $\Gamma$は$␣ \in \Gamma$かつ$\Sigma \subseteq \Gamma$であるようなテープ・アルファベット

				\item $\delta : Q \times \Gamma \to Q \times \Gamma \times \{ L, R \}$は遷移関数

				\item $q_0 \in Q$は開始状態

				\item $q_{\mathrm{accept}} \in Q$は受理状態

				\item $q_{\mathrm{reject}} \in Q$は$q_{\mathrm{reject}} \neq q_{\mathrm{accept}}$であるような拒否状態
			\end{enumerate}

		\end{dfn}

        また，your headsと複数形になっていることは，単純にyour eyesから複数形を引き継いでいるだけとも考えられるが，読者それぞれの頭で固有の物語を考えてほしいという意図の表れとも読める．

        \subsection{エピグラフ}\label{ed}

            \begin{quote}
                「維特因図霊問狗子還有言語也無維云無」
            \end{quote}

            南宋の僧侶，無門慧開による『無門関』に収められた48の公案のうち，最初の「狗子仏性」\footnote{夏目漱石『夢十夜』の第二夜にも登場する．}をもじったもの\footnote{「趙州和尚，因僧問，狗子還有仏性也無．州云，無．」「趙州和尚，因みに僧問う，「狗子に還って仏性有りや」．州云く，「無」．」\cite{mumonkan}}．越州和尚がヴィトゲンシュタイン\footnote{維特はヴィトゲンシュタインの中国語表記（維特根斯坦）の一部．ゲーテ『若きウェルテルの悩み』のウェルテルも中国語では維特と表記されるが，チューリングと並べる人物としてはヴィトゲンシュタインに比べると不適．ただし，『若きウェルテルの悩み』も恋愛に苦しむ人物を描いた小説であることから，一定程度注意を払うべきだろう．}に，僧がチューリング\footnote{図霊はチューリングの中国語表記．}に，仏性が言語に置き換わっている．

            ヴィトゲンシュタインが導入されていることについて．ヴィトゲンシュタイン『哲学探究』において，犬を例に挙げて議論を行なう箇所がある．これは，犬は痛いふりが出来ない，なぜなら犬の生活は単純すぎて痛いふりをする複雑な機能がないからだ，というもの．円城塔が頻繁に用いる，ある者がある事物を想像出来ないとき，その者にとってその事物は存在するのだろうか，という議論にも通じる．また，この議論はチューリングテストにも似ている．

            さらに，『哲学探究』が後期ヴィトゲンシュタインの著作であることに注意したい．つまり，円城塔は，ここからは言語ゲームを展開するぞ，という意図の開陳としてヴィトゲンシュタインを提示している可能性がある．一方で，前期ヴィトゲンシュタインは強力な独我論を展開した哲学者であり，独我論者であるレフラーを象徴する人物とも言える．つまり，「Boy's Surface」から「Your Heads Only」を繋ぐ象徴として導入している可能性も指摘される．

            チューリングが導入されていることについて．先述の通り，ヴィトゲンシュタインの議論がチューリングテストに似ていることからの連想であるとともに，本作にチューリングが登場するための布石でもある．チューリングとヴィトゲンシュタインは同時代の人物であり，ヴィトゲンシュタインの講義にチューリングが出席していたこともあった．%ウィトゲンシュタインvsチューリング, 勁草書房, 2023
            また，チューリングは同性愛者であることを理由に迫害されたことが知られており，ヴィトゲンシュタインについても同性愛傾向が指摘されている．

            「狗子仏性」について．無門は，禅を通じて悟りを得たいのであれば，まず意識を徹底的になくすことが必要であるとする．先人たちも通った最初の関門がこの無であるとし，全身を疑いの塊とし，無になること（意識をなくすこと）を強く求める．昼夜の別なくこの問題に取り組み，虚無や有無といったようなものを考えることから離れ，これまでに身に着けてきたものを削ぎ落していき，自然に無に行き着くのだという\footnote{この姿勢には，数学者・理論物理学者が目標とする姿勢と共通するものがある．数学者の佐藤幹夫は，次のように述べた．「朝起きた時に，きょうも一日数学をやるぞと思ってるようでは，とてもものにならない．数学を考えながら，いつのまにか眠り，朝，目が覚めたときは既に数学の世界に入っていなければならない．」\cite{kimura}}\cite{mumonkan}．

            趙州は，ここでは犬に仏性はないとしているが，別の機会にはあると答えたという．このように，趙州が犬の仏性について矛盾する回答を同等に与えたことは，後世において有無の二元論を超えた“無”を示したと解釈された．このことから，狗子仏性は，二元論のような論理が確立する以前の，体験の世界に導くための公案であるとされるようになった\cite{mumonkan}．

        \subsection{p.136, l.1}

            \begin{quote}
                「私にはそんなところに〜」
            \end{quote}

            1999の語り手は“私”．円城塔によれば，本作は“僕”，“私”，無人称の語り手が3回ずつ登場し，物語とは何ら関係のない三目並べを行っているという\cite{sri}．

        \subsection{p.136, l.6-9}

            \begin{quote}
                「その断章が笑えるものであったなら，一を記す．〜」
            \end{quote}

            本作を読み進めるためのルール．テクストがそう読まれることを望んでいるということで，実際に手を動かして読み進めるべきだ．

            なお，テクスト自身がそのテクストを読むためのルールを指定してくるテクストとして，フリオ・コルタサル『石蹴り遊び』が挙げられる．

        \subsection{p.137, l.6}

            \begin{quote}
                「私はそんな種類の兵器ではない．」
            \end{quote}

            ここでは，何度読んでも同じ情動を引き起こさせる文字列を“兵器”と表現しているものと推察される．読書とはテクストと読者の相互作用によって記述され，その相互作用によって読者は不可逆に変化する．たとえ同じテクストを再び読んだとしても，読者は既に相互作用によって変化しているので，読者が同じ感想を得るかどうかは不明である．逆に，何度読んでも同じ情動を引き起こさせる文字列があるとすれば，それは確かに何らかの兵器である，という感が強い．

            ここで，円城塔作品が常々「わからない」というお決まりの評価を下されていることを考えると，どの作品を何度読んでも「わからない」と評価されてしまう円城塔の作品は，一種の兵器であると言える．読書をテクストと読者の相互作用で記述しようという姿勢は，「Boy's Surface」「ドルトンの印象法則」でも扱われる主題である．

            円城塔作品を一種の兵器に例えることについて，以下のような事実を紹介したい．宮本輝は，2011年上半期の芥川賞選考会において，「これはペンです」について，次のように評している\cite{bunshun1}．

            \begin{quote}
                私は，書き出しのたったの十ページを読んだだけで眠くなり，全篇読了は難行苦行だった．要するに，つまらなかったのだ．
            \end{quote}

            これに対して，島田雅彦は，2011年下半期の芥川賞選考会において，「道化師の蝶」について，次のように述べている\cite{bunshun2}．

            \begin{quote}
                二回読んで，二回とも眠くなるなら，睡眠薬の代わりにもなる．
            \end{quote}

            島田雅彦の論は回を跨いだ宮本輝への反論であるように思われる．しかも本作が睡眠薬代わりとなる一種の兵器であると述べているように捉えられる．ただし，宮本，島田の論はいずれも本作発表後のものであることに注意せよ．

            逆に，円城塔は，自身の作品が決まり文句のように「わからない」「読むと眠くなる」と繰り返し聞かされ，開き直っているように読める．

            「パリンプセストあるいは重ね書きされた八つの物語」の「涙方程式始末」に登場する，目にするだけで涙を流してしまう“涙多様体”もこの兵器の一種であろう．何度読んでも読者に対して同じように作用する文字列というモチーフは，初期円城塔が好んで用いるもののひとつである．

        \subsection{p.139, l.1}\label{rule90}

            \begin{quote}
                「ルール90」
            \end{quote}

            1次元3近傍セルオートマトンのルール90のこと．後続のルール90に関する説明はいずれも正しい．ここで，ルール90とは，以下のように定められる．

            \begin{dfn}{ルール90}{}
                1次元3近傍セルオートマトンのルール90は以下の通り定められる．なお，上段が現在の状態，下段が次の状態である．
                \begin{align*}
                    \blacksquare &\blacksquare \blacksquare & \blacksquare &\blacksquare \Box & \blacksquare &\Box \blacksquare & \blacksquare &\Box \Box & \Box &\blacksquare \blacksquare & \Box &\blacksquare \Box & \Box &\Box \blacksquare & \Box &\Box \Box\\ &\Box & &\blacksquare & &\Box & &\blacksquare & &\blacksquare & &\Box & &\blacksquare & &\Box
                \end{align*}
            \end{dfn}

            ルール90との名前は，$\Box$を0，$\blacksquare$を1というように2進数で表記し，上記のように現在の状態について降順に並べたとき，次の状態が01011010，つまり2進数で90と書けることによる．このような命名法を，ウルフラム・コードという\footnote{ウルフラムの著書``A New Kind of Science''\cite{nks}で提案されたことに由来する．}．

            また，十分大きな空間に1つだけ$\blacksquare$が置かれている場合のルール90の時間発展は，以下のようにシェルピンスキーのギャスケットを生成する．なお，シェルピンスキーのギャスケットとは，正三角形の各辺の中点を結び，その新たに出来た正三角形を取り除く操作を再帰的に実行することで得られるフラクタル図形である．

            \begin{exm}{ルール90の時間発展}{}
                \begin{gather*}
                    \Box \Box \Box \Box \Box \Box \Box \Box \Box \blacksquare \Box \Box \Box \Box \Box \Box \Box \Box \Box\\[-0.7em]
                    \Box \Box \Box \Box \Box \Box \Box \Box \blacksquare \Box \blacksquare \Box \Box \Box \Box \Box \Box \Box \Box\\[-0.7em]
                    \Box \Box \Box \Box \Box \Box \Box \blacksquare \Box \Box \Box \blacksquare \Box \Box \Box \Box \Box \Box \Box\\[-0.7em]
                    \Box \Box \Box \Box \Box \Box \blacksquare \Box \blacksquare \Box \blacksquare \Box \blacksquare \Box \Box \Box \Box \Box \Box\\[-0.7em]
                    \Box \Box \Box \Box \Box \blacksquare \Box \Box \Box \Box \Box \Box \Box \blacksquare \Box \Box \Box \Box \Box \\[-0.7em]
                    \Box \Box \Box \Box \blacksquare \Box \blacksquare \Box \Box \Box \Box \Box \blacksquare \Box \blacksquare \Box \Box \Box \Box \\[-0.7em]
                    \Box \Box \Box \blacksquare \Box \Box \Box \blacksquare \Box \Box \Box \blacksquare \Box \Box \Box \blacksquare \Box \Box \Box \\[-0.7em]
                    \Box \Box \blacksquare \Box \blacksquare \Box \blacksquare \Box \blacksquare \Box \blacksquare \Box \blacksquare \Box \blacksquare \Box \blacksquare \Box \Box \\[-0.7em]
                    \Box \blacksquare \Box \Box \Box \Box \Box \Box \Box \Box \Box \Box \Box \Box \Box \Box \Box \blacksquare \Box \\[-0.7em]
                    \blacksquare \Box \blacksquare \Box \Box \Box \Box \Box \Box \Box \Box \Box \Box \Box \Box \Box \blacksquare \Box \blacksquare \\
                    \cdot \\[-0.7em] \cdot \\[-0.7em] \cdot
                \end{gather*}
            \end{exm}

        \subsection{p.139, l.4-5}

            \begin{quote}
                「今でも彼らにルール90と話しかけると，自然と何かが笑み零れてくる．」
            \end{quote}

            “ルール90”という言葉が，一部の人間に対しては兵器として作用することを述べている．確証はないが，概ね事実である．私自身，ルール90と言われたら，少なくとも嫌な気持ちにはならないし，生成されるシェルピンスキーのギャスケットを想像してしまって自然と面白く思われてしまう．

        \subsection{p.139, l.14-16}

            \begin{quote}
                「日本語としては〜」
            \end{quote}

            日常語とはまったく異なる，不思議な語感をもつ言葉を日常的に操る人間は確かに存在していて，しかも普通に日常生活を送っている．既にお気づきの通り，あなたは今，その奇妙な日本語が用いられた文章を読んでいる．

            確かに存在しているのだが，透明化されて社会にはいないことになっているいわゆる理系\footnote{数学が得意であるような人間，あるいは数学が苦手ではない人間と言い換えてもよい．}に関する話題は，のちのち『土人形と動死体』で主題として扱われる．

        \subsection{p.140, l.2-3}

            \begin{quote}
                「私とそんな人々は，敵対関係と呼ばれる状況へと吹き寄せられつつあるのだが，」
            \end{quote}

            “私”を円城塔と仮定すれば，これは円城塔が小説家としてのキャリアを確立しつつあることを述べている．小説家およびその生成物である小説はいわゆる文系の領分であり，これはいわゆる理系とは相反するもので，しかも敵対するものと考えられている．このことを茶化した記述だろう．

        \subsection{p.140, l.14}

            \begin{quote}
                「ルール110」
            \end{quote}

            既に\ref{rule90}節で述べた通り，ルール110もまた1次元3近傍セルオートマトンのルールである．念のため，ルール110は以下の通り定められる．

            \begin{dfn}{ルール110}{}
                1次元3近傍セルオートマトンのルール110は以下の通り定められる．なお，上段が現在の状態，下段が次の状態である．
                \begin{align*}
                    \blacksquare &\blacksquare \blacksquare & \blacksquare &\blacksquare \Box & \blacksquare &\Box \blacksquare & \blacksquare &\Box \Box & \Box &\blacksquare \blacksquare & \Box &\blacksquare \Box & \Box &\Box \blacksquare & \Box &\Box \Box\\ &\Box & &\blacksquare & &\blacksquare & &\Box & &\blacksquare & &\blacksquare & &\blacksquare & &\Box
                \end{align*}
            \end{dfn}

            ルール110は十分複雑であり，ウルフラムによるオートマトンの分類ではクラス4に分類される．実際，ルール110はチューリング完全であることが知られている\cite{cook}．

        \subsection{p.141, l.1-3}

            \begin{quote}
                「非道く当たり前のことではあるのだが，〜」
            \end{quote}

            先述の，相互作用としての読書と，同じ感想を何度も得ることが出来るだろうか，という議論そのもの．「ドルトンの印象法則」の主題．

        \subsection{p.141, l.10-11}

            \begin{quote}
                「世の中にはいわゆる叙述トリックというものが全く通用しない人たちがおり，〜」
            \end{quote}

            自閉スペクトラム症者のことを指すか．自閉スペクトラム症者に見られる特徴の1つとして，比喩表現を比喩であるとわからないこと，つまり文章を文字通り読んでしまうことが知られている（ウイングの三つ組\footnote{医師で自閉スペクトラム症児の母であるウイングによって報告\cite{wing}された，自閉スペクトラム症者の特徴のこと．空気が読めず，他人に興味を持たないという\textbf{社会性の障害}，言語獲得期に鸚鵡返しが見られる，自然言語を字義通りにしか解釈しないことによって比喩や冗談が通じないという\textbf{コミュニケーションの障害}，目の前にないものを想像出来ないという\textbf{想像力の障害}の三つを指す\cite{hongo}．円城塔作品には，これらの特徴を顕著に示す人物が多数登場することに特に注意したい．}）．ここに続けて“文字通り読め”ということを繰り返し述べているところから，この見立ては十分に正当化される．なお，理学・工学の研究者・技術者は自閉的傾向が強いことが知られており，逆に自閉的傾向が理学・工学への適性の高さに繋がっているのではないかとの議論もある\cite{buchen}．

            また，ディラックはまさに小説を文字通りに読んだという有名な逸話がある．ディラックは，ドストエフスキー『罪と罰』を読み，「これは良い作品であるが，ある章で著者は誤りをおかしている．同じ日に二度太陽が昇るように記述している」と語った\cite{dirac}p.59．
            円城塔は，小説は人間の認知能力に依存するメディアであり，人間の認知能力が変化すれば，小説もまた自然と変化するだろうと繰り返し主張している\cite{interview1,interview2}．この主張に加えて，自閉スペクトラム症の認知能力が定型発達者の認知能力と異なると仮定すれば，自閉スペクトラム症者の創作した小説は定型発達者からすれば奇妙なものである可能性があり，そもそも自閉スペクトラム症者が記述する世界自体が定型発達者からすれば奇妙な世界である可能性がある，ということが導かれる\footnote{小説の創作ではなく読解に関してではあるが，自閉スペクトラム症者は定型発達者とは異なる視座を持っていることが指摘されている\cite{savarese}．}．円城塔作品が「難解」「わからない」と評されるのは，このような，通常とはまったく異なる感覚\footnote{ここで，円城塔本人の発達特性については論じていないことに十二分に注意せよ．“通常とはまったく異なる感覚”として最初に検討すべきものは，円城塔が暗黙裡に多用する相対論・量子論及び内在物理学である．}によって創作されていることが原因である可能性が指摘される．

        \subsection{p.142, l.7}

            \begin{quote}
                「注文の多い料理店」
            \end{quote}

            無論，宮沢賢治「注文の多い料理店」．

        \subsection{p.142, l.11-13}

            \begin{quote}
                「一つ．手を叩く音が聞こえたら，何をおいても笑い出す．／二つ．その断章が全くおかしくなかったならば，一つ前の断章に戻って読み返す．／三つ．こんな話につきあうのが嫌になったら，この文章を読み始める前にしていた作業へ真っ直ぐ戻る．」
            \end{quote}

            本作を読むための簡単なルール．文字通り守るべきだろう．

        \subsection{p.143, l.6-7}

            \begin{quote}
                「これはいわゆる免責事項ということになる．」
            \end{quote}

            佐々木敦が『あなたは今，この文章を読んでいる．』\cite{sasaki}で定義した，パラテクストの代表例はREADME\footnote{ソフトウェアに内包される形で配布される，そのソフトウェアの注意事項などを記した取扱説明書的な文書．使用者に最初に読まれるために作成・配置される文書であり，使用者が最初に必ず読む文書である．GitHubには，リポジトリ直下にREADMEが存在する場合はREADMEをリポジトリの説明として表示する機能がある．}であった．免責事項も，誰かに読まれることを前提として書かれ，そして実際に読まれるという点でパラテクストと言える．パラテクストを意図的に作中で提示していることに注意したい．

        \subsection{p.143, l.10}

            \begin{quote}
                「私はここで唐突に一つ手を打ち鳴らす．」
            \end{quote}

            手を叩く音が聞こえたので，我々は笑い出さなくてはならず，しかも，笑うならばおかしいので，我々は次の断章に進まなければならない．最初の断章は$\Box$としなければならない．%これは，心理学における抹消起源説と中枢起源説の紹介．二要因説の吊り橋効果の話が本作の論理展開に酷似していることを指摘．
            %https://bsd.neuroinf.jp/w/index.php?title=%E6%83%85%E5%8B%95&mobileaction=toggle_view_desktop

            円城塔は，ある因果関係について，その因果関係を1対1対応にすり替え，本来の結果から原因を得るような論理展開を多用する．「文字渦」における，漢字とその指示対象との1対1対応にその好例が見える．

        \subsection{p.143, l.17-p.144, l.1}

            \begin{quote}
                「今あなたの耳に，〜」
            \end{quote}

            先ほどの手を叩く音が音に聞こえなかった場合への対策．読者は強制的に次の断章に進む．

        \subsection{p.144, l.6}

            \begin{quote}
                「僕の中の記憶の彼女は，〜」
            \end{quote}

            2000の語り手は“僕”．

        \subsection{p.144, l.6-7}

            \begin{quote}
                「あれは僕が博士課程を修了してなんとか一度目のポスト・ドクターの口にありついていた頃の出来事〜」
            \end{quote}

            ポスト・ドクターとは，博士課程を修了（博士号を取得）するか博士課程を研究指導認定退学\footnote{かつて単位取得退学と言っていたもの．}した後の進路のひとつで，大学・研究所等での有期雇用のこと．略してポスドク．

            民間企業への就職や大学・研究所等での無期雇用（テニュア）が確保できなかったもののアカデミアで生きていこうとする場合はポスドクを選ぶことになる．長年の勉学と研究の果てにある最終到達点ではあるのだが，ポスドクの給与は極めて低く，また有期雇用契約であるため非常に不安定な身分となっている．円城塔が日本物理学会誌に寄稿したエッセイ「ポスドクからポストポスドクへ」\cite{posdoc}は，ポスドクの悲惨な実態を記している\footnote{なお，同作は法螺吹きを専門とする円城塔によるフィクションであり，こんな地獄がわが国に顕現しているはずがない．}．

            また，円城塔が博士課程を修了し，最初のポスドクについたのが2000年であり，文中の記述と一致していることに注意せよ．さらに，各章の語り手が同一ではない可能性にも注意したい．

        \subsection{p.145, l.7}

            \begin{quote}
                「僕自身はおそらく，あらゆる対象への好悪が薄く，〜」
            \end{quote}

            他人への興味が薄く，対人関係に関する語彙が極端に少ないところから，この語り手もまた自閉スペクトラム症者である可能性が指摘される．

        \subsection{p.146, l.4-5}

            \begin{quote}
                「周囲からは半ばゲイ扱いをされたりもした〜」
            \end{quote}

            語り手の性別が確定されているのは初期作品特有．『文字渦』以降，登場人物の性別があからさまに判断出来る作品はほぼ存在しなくなる．

        \subsection{p.146, l.11-12}

            \begin{quote}
                「生まれた土地から離れて学んだことは，人間以上という区分が存在すること．」
            \end{quote}

            “人間以上”はシオドア・スタージョン『人間以上』を意識したものか．

        \subsection{p.147, l.3}

            \begin{quote}
                「通常，もし数学ができたならに続けられるのは，割り算が早くて〜」
            \end{quote}

            いわゆる理系の人間，特に数学科・物理学科の人間にしばしば浴びせられる言葉の代表格．“それは算数や算術であって数学ではない”などという反論が聞き入れられないことは明白であり，どう転んでも場の空気を好転させることはないので押し黙るのが最善策となる．

            数学者・物理学者の大半はそのような計算は苦手であり，算術が得意なのは工学者や天文学者ではなかろうか．なお，数論を専門とする数学者や理論物理学者の一部には算術が異常に得意な者がいる．歴史上の科学者では，ガウス，フォン・ノイマン，フェルミ，ファインマンが得意とする一方で，ポワンカレ，ヒルベルト，グロタンディーク，アインシュタインは不得意だった（少なくとも大得意ではなかった）ことが知られている．

        \subsection{p.147, l.3-5}

            \begin{quote}
                「お会計向きの足し算割り算の(中略)四千年前のパリコレで発表済みであることが知られている．」
            \end{quote}

            “〜が〜であることが知られている”は数学・物理学の教科書で頻出するフレーズ\footnote{既に本書でも何度か出現している．}．その分野で広く知られている常識を学習者に対して提示するためのお決まりの文句．証明を省略するときの常套句の一つでもある．

            類例としては“証明は読者の演習問題とする”，“賢明な読者諸君にはおかれては既にお気づきの通り”．

        \subsection{p.147, l.8-10}

            \begin{quote}
                「物理屋と数学者の間の溝は冥くて深く，物理屋と生物学者の間の距離は更に広い．〜」
            \end{quote}

            物理屋とは物理学者が物理学者のことを指すときに使う，やや自虐の入った自称・他称．さらに細かく，実験屋と理論屋という言葉もある．

            元々，仮説（理論）の構築とその実証実験は一人の科学者が一身でこなしていた．20世紀以降，数学・物理学の急速な発展によって，理論に用いる数学と実験に用いる工作技術が極端に複雑になり，まず数学と物理学が分離された\footnote{正規教育をほとんど受けていないファラデーを例外として，19世紀末の解析力学・電磁気学・熱力学の頃までは数学者と物理学者はほとんど区別されなかった．ガウス，コーシー，ストークス，リウヴィル，カラテオドリ，ケルヴィン卿などがその代表格である．19世紀末から20世紀初頭の数理論理学・量子論・相対論の急速な発展によって，数学者と物理学者は急速に分化独立した．以降，ヒルベルト，フォン・ノイマン，ヴァイル，ウィッテンなどを例外として，数学者と物理学者を兼ねることは非常に困難となっている．}．物理学の中でも理論と実験の両立が不可能となって，実験の専門家（実験物理学者，実験屋）と理論の専門家（理論物理学者，理論屋）に別れた．理論と実験の双方で一流の業績を挙げた最後の物理学者が，エンリコ・フェルミであった．

            大学の学部では，物理学科・数学科・生物学科は理学部や理工学部にまとめられているが，その内部では互いに互いを遠い存在だと思い合っている．物理屋は数学者の数学に対して，物理学的には瑣末なところに異常に拘ってどうしようもないというような感想を抱きがち．逆に，数学者は物理屋の数学に対して，雑すぎて話にならないと思っているはず．物理屋は生物学に対して，分子生物学や生物物理学といった形で侵略を試み，相当の成功をおさめている．これに対して生物学者からの抵抗は大きかったが，現在では分子生物学の成果は当然のものとみなされているように思う．なお，円城塔の指導教官である金子邦彦の専門は生物物理学・理論生物学であり，研究室も生物物理学系である．

            ところで，自然の階層性を考慮すると，数学，物理学，化学，生物学という順番になるはずだが，本作では化学が省略されている．これは，金森修『サイエンス・ウォーズ』で次のような例え話があることに由来していると推察される\cite{kanamori}p.96．

            \begin{quote}
                例えば海産無脊椎動物の生態を調査している生物学者と理論物理学者との間の距離は、物理学者とガリソン\footnote{科学哲学者のピーター・ギャリソンのこと．客観性という概念が近代以降急速に構築された新しい概念であることを論じた『客観性』（ロレイン・ダストンとの共著）\cite{objectibity}が有名．}という科学論者との距離よりもはるかに大きいに違いない。
            \end{quote}

            また，生物学者は物理学者を，物理学者は数学者を魔法使いだと思っているというのもその通り．生物学，物理学だけで上手くいかなかったときは魔法使いに相談に行く．生物学ではシュレーディンガーに端を発する分子生物学\footnote{ただし，分子生物学については，（物理学を半ば追放された）シュレーディンガーが生物学に殴り込んで行って主張したものなので，厳密には事例としては不適．}が，物理学ではリーマン幾何学を用いた一般相対論，群論を用いた量子力学がその実例である．

        \subsection{p.148, l.6-7}

            \begin{quote}
                「薄茶色をした薄い殻に囲まれて，〜」
            \end{quote}

            マジで人間の外見描写なのか？

        \subsection{p.150, l.6-p.151, l.16}

            \begin{quote}
                「土斑猫．〜」
            \end{quote}

            ツチハンミョウの奇妙極まる生態に関する説明．作中の説明は正しい．ツチハンミョウの生態は，ジャン-アンリ・ファーブル『昆虫記』\cite{fabre}で有名．

        \subsection{p.152, l.17}

            \begin{quote}
                「雪片曲線．有限の面積を囲む，至る所微分不可能な線．」
            \end{quote}

            コッホ雪片とも．線分を3等分し，中央の部分線分を他の部分線分と同じ長さの正三角形の2辺と置き換える操作を繰り返すことで得られるフラクタル曲線をコッホ曲線といい，正三角形の各辺に対してこの操作を行う，つまりコッホ曲線を3つ組み合わせることによって得られる閉曲線をコッホ雪片という．コッホ雪片に囲まれた内部は有限の面積を持ち，またコッホ曲線はいたるところ微分不可能である．ここで，“いたるところ微分不可能”は厳密な定義語である．

            \begin{dfn}{いたるところ微分可能，いたるところ微分不可能}{}
                関数$f$の定義域を$I$と書く．$f$が\textbf{いたるところ微分可能}であるとは，任意の$x \in I$について$f$が微分可能であることである．また，$f$が\textbf{いたるところ微分不可能}であるとは，任意の$x \in I$について$f$が微分不可能であることである．
            \end{dfn}

            収束・連続・微分可能といった解析学の中核となる概念は，コーシーとヴァイエルシュトラスによって，19世紀をかけて$\varepsilon$--$\delta$論法によって厳密化された．コーシーは$\varepsilon$--$\delta$論法の導入に大きく貢献しているが，コーシー自身の議論には不徹底なところがあり，しばしば誤謬が見られる\index{これは飯田隆『言語哲学大全1』\cite{iida1}で真っ先に指摘されている．}．コーシーの誤謬の具体例として，アンペールの定理を認めていたことが挙げられる\cite{nakane}．ここで，アンペールの定理は以下の通り．

            \begin{clm}{アンペールの定理}{}
                任意の連続関数は，有限個の特異点を除いて微分可能である．
            \end{clm}

            “領域$I$で関数$f$が常に微分可能\,\to\,$I$で$f$が連続”は自明だが，その逆，つまり“$I$で$f$が連続\,\to\,$I$で$f$は常に微分可能”（アンペールの定理より強い主張）は非自明であり，しかも成り立たないことが知られている．アンペールの定理はアンペール，コーシー，ベルトラン，ガウス，ディリクレといった優秀な数学者たちの注意をすり抜け，長期間にわたって生き残っていた．これに対して，ヴァイエルシュトラスは，アンペールの定理の反例として，いたるところ微分不可能な関数であるヴァイエルシュトラス関数\footnote{このような，満たしてほしい性質を満たさない，つまり振る舞いの悪い性質をもつ関数を病的(ill-behaved)な関数という．ある対象を病的であると判断する基準は数学者の感覚に依存している．ただし，ある性質が病的と判断された場合，病的な例を正しく取り込む形で理論が修正される．ヴァイエルシュトラス関数などの病的な関数の存在を踏まえて，アンペールの定理は破棄され，微分可能・連続などの概念の厳密化を経て，病的な関数は解析学でも正しく扱える対象となっていった．このような営みは，対象と理論の共進化であると言えるかもしれない．}を実際に構成し，アンペールの定理を棄却した．ヴァイエルシュトラス関数とコッホ曲線は同じいたるところ微分不可能な関数であり，数学的に正しく扱うためには$\varepsilon$--$\delta$論法が不可欠である．

        \subsection{p.153, l.11}

            \begin{quote}
                「ここに，一つの取り残された二十世紀が挟み込まれて，〜」
            \end{quote}

            「Gernsback Intersection」の二十世紀のこと．また，2001の語り手は無人称．

        \subsection{p.154, l.5}

            \begin{quote}
                「ルシフェラーゼ」
            \end{quote}

            緑色蛍光タンパク質(Greem Fluorecent Protein, GFP)のこと．GFPは，「緑字」「捧ぐ緑」にも登場する．

            円城塔作品に緑色蛍光タンパク質が登場する理由について．生物を何らかのダイナミクスとして記述することを考える．従来，このようなダイナミクスとしては，生物の個体数のようなマクロな測定値しか扱えなかった\footnote{生物物理学の古典的研究の最も有名な例として，生物の個体数と資源に関するロトカ\hyphen ボルテラ方程式が挙げられる．これが直接の題材となっているのが「内在天文学」．}．

            しかし，下村脩によってGFPを利用する手法が開発された\footnote{この業績によって，下村は2008年にノーベル化学賞を受賞した．なお，下村と私自身には親戚関係はない．}ことで，ミクロな測定値も扱えるようになった．例えば，あるタンパク質に注目したいとき，そのタンパク質とGFPが結合したタンパク質を合成するように遺伝子改変を行えば，GFPによる蛍光の量を注目したいタンパク質の量とみなすことが出来，細胞等を生かしたままミクロな測定を行うことが可能になった\cite{cell}．

            GFPが円城塔作品において何らかの情報を示すマーカーの代表例として頻出するのは，上記のようにGFPが複雑系研究を飛躍的に発展させたことを背景としている．

        \subsection{p.154, l.6}

            \begin{quote}
                「キルリアン写真に念写に液体からの伝言」
            \end{quote}

            疑似科学3連発．初期円城塔には，疑似科学を馬鹿にした記述が散見される．

            キルリアン写真とは，フィルムや印画紙に生体の一部を押し付け，高電圧を印加しながら撮影することによって得られる写真のこと．キルリアン写真では生体の周囲に特徴的な影が生じ，これを生体の“オーラ”を捉えたものであるとする．名前はこれを1940年頃に発見したソ連のキルリアン夫妻に因む．なお，キルリアン写真は湿気に超高電圧を印加したことによるプラズマ放電を捉えたものであり，生体の“オーラ”の撮影ではない\cite{boyers}．

            念写とは，頭で思い浮かべた画像をフィルム等に写し出すこと，またそれによって得られた写真のこと．心理学者で東京帝大助教授の福来友吉によって報告された．福来は著書『透視と念写』\cite{fukurai}において，透視も念写も事実であると主張した．これに対して，東京帝大物理学科の教官・学生を中心とした人物らによる反論書\cite{fuji}では徹底的に批判されている．福来は文官分限令に基づき休職処分となり，最終的に東京帝大・学会から追放されている\footnote{分限令の理由は定かではないが，疑似科学を繰り返し主張したことによって東京帝大の権威を傷つけたことによると見るのが妥当だろう．}．

            「液体からの伝言」は，江本勝の著書『水からの伝言』\cite{emoto}を指す．江本は，同書において，「ありがとう」という言葉を聞かせた水を冷却させて得られる氷の結晶は綺麗な結晶となるが，「ばかやろう」と聞かせた水を冷却させて得られる氷の結晶は崩れた結晶となる，と主張する．これはもちろん疑似科学である\footnote{\cite{amabane}に見える通り，『水からの伝言』の批判者として，物理学者で複雑系の研究者である菊池誠，物理学者で複雑系の批判者である田崎晴明が有名．菊池はSFファンダムで長く活動している人物であり，円城塔と直接の知り合いであることにも注意．田崎もまたSFファンである（田崎の英名Hal Tasakiはクラーク『2001年宇宙の旅』のHALに由来する）．また，菊池は統計力学の応用である複雑系の研究者であり，田崎は統計力学の理論を専門としていて学問的な基礎はまったく同根であること，田崎の複雑系批判は主に東大の複雑系，つまり金子邦彦・池上高志らに向けられたものであることにも注意．そもそも，物理学者たちによる議論は，いかに熾烈に見えたとしても，物理学という舞台で闘っている限りは禍根を残すようなものではないことにも注意したい．}が，少なくとも2000年代前半においては，一部の小学校・中学校における道徳・総合学習などに取り入れられていた\cite{amabane}．

        \subsection{p.154, l.9}

            \begin{quote}
                「月のチーズに夏の欠片」
            \end{quote}

            前者は実在する商品名，後者はV6の楽曲名に同名のものが存在するが．真にこれを意図したものかは不明．

        \subsection{p.155, l.16-17}

            \begin{quote}
                「とある解説に従うならば，二十世紀は第四コーナーを曲がり損ねて，盛大にオーバーランしてクラッシュした．」
            \end{quote}

            「Gernsback Intersection」の物語そのもの．“とある解説”が何を指すかは不明だが，20世紀の第4四半期に“大きな物語”は終焉したとするポストモダン思想が流行したことから，この“大きな物語”の終焉という言説が“とある解説”であると解釈するのが妥当だろう．

            藤田直哉\cite{fujita}が指摘するように，円城塔は，露骨にポストモダン的解釈を誘うような書き振りをし，実際ポストモダン的に解釈出来るようにテクストを構成しつつも，最後まで古典的な操作の徹底を求める．この箇所についても，ポストモダン的な解釈を誘っていると解釈していいだろう．

        \subsection{p.157, l.10-11}

            \begin{quote}
                「東には仏教的エデンがある」
            \end{quote}

            仏教における東方浄土のことか．「Gernsback Intersection」に西方浄土についての言及があったことに注意したい．

        \subsection{p.157, l.12}

            \begin{quote}
                「我々はエデンの東にいる」
            \end{quote}

            ジョン・スタインベックの長篇小説『エデンの東』．

        \subsection{p.157, l.15-17}

            \begin{quote}
                「ただ瓶だけが流れ着き続ける．やがて臨界数量に達することを得た瓶たちは，〜」
            \end{quote}

            ウルフラムの提唱によって，十分大きな瓶の集まりはチューリング完全となり，生命となって言葉を語り始める．ここで，ウルフラムの提唱は以下の通り．

            \begin{ths}{ウルフラムの提唱}{}
                十分複雑なものはチューリング完全であり，それは生命である．
            \end{ths}

        \subsection{p.159, l.1-2}

            \begin{quote}
                「ロリータ，我が命の光，〜」
            \end{quote}

            途中までは若島正訳のウラジーミル・ナボコフ『ロリータ』\cite{lolita}の冒頭の引用だが，最後に『三百六十五歩のマーチ』をもじったギャグになる．

        \subsection{p.159, l.4-6}

            \begin{quote}
                「僕は三キロほどはあろうかというビニール袋と，〜」
            \end{quote}

            2002の語り手は“僕”．

        \subsection{p.159, l.14-p.160, l.2}

            \begin{quote}
                「当時はまだ助手と〜」
            \end{quote}

            アカデミアの貧困化への皮肉．初期作品の随所に同様の記述が見える．

        \subsection{p.160, l.15}

            \begin{quote}
                「チューリングにお供えするなら林檎だろう」
            \end{quote}

            チューリングが青酸を塗布した林檎を食べて自殺した事実から．やや不謹慎．

        \subsection{p.162, l.13}

            \begin{quote}
                「局所的には最適解へ落ち込んでいる．」
            \end{quote}

            最適化問題を解くときのあるある．複雑系におけるアトラクタへの吸引もこの類似．

        \subsection{p.163, l.11}

            \begin{quote}
                「槍形吸虫」
            \end{quote}

            カタツムリなどの陸貝の体内で繁殖し，アリを経て大型哺乳類に寄生する寄生虫の一種\cite{waki}．作中の説明はすべて正しい．

            作中において，槍形吸虫の奇妙な生活史を機械論的に紹介しているのは，「Boy's Surface」「Goldberg Invariant」「Gernsback Intersection」で繰り返し示される，“可能なことはいつか必ず実現される”という物理学者が暗黙理に抱いている自然観を強調する狙いがある．カタツムリ，アリ，ヒツジはそれぞれある法則によって行動しているので，それらから成る系においては，一見奇妙ではあるが成り立ちうる生活史があり，その生活史を満たす生物がいつか必ず実現される．

            また，いかに他人からは奇妙に見えても，本作を通じて描かれるようないわゆる理系の人々がこの社会に実在しているという主張の隠喩としても読める．

        \subsection{p.165, l.7-8}

            \begin{quote}
                「彼の名を冠した，名誉に満ち溢れた賞」
            \end{quote}

            チューリング賞．計算機科学において，最も名誉ある賞．

        \subsection{p.165, l.16}

            \begin{quote}
                「ハンバート・ハンバートのどうしようもない性癖」
            \end{quote}

            ハンバート・ハンバートはナボコフ『ロリータ』の語り手．少女ロリータへの執着から小児性愛者とされる．

        \subsection{p.167, l.1}

            \begin{quote}
                「ベニクラゲの特性」
            \end{quote}

            ベニクラゲは不老不死であることが知られている．

        \subsection{p.168, l.8-9}\label{malthus}

            \begin{quote}
                「この年は，全弾のいちいちが命中して，十月十日後の人類滅亡が〜」
            \end{quote}

            受精確率が極端に増大し，急激な人口爆発が発生するということ．マルサスの罠へと連想が続いていく．また，2003の語り手は無人称．

        \subsection{p.169, l.17-p.170, l.1}

            \begin{quote}
                「数を順に数えていたらいつの間にか無限を数え過ぎていた〜」
            \end{quote}

            既に何度か登場している，数え上げによる無限の定義を前提としたギャグ．自然言語における“無限”が数学的にどのような意味で用いられているのかを確定することは難しく，昔から多くの不毛な議論があった．詳細はカントール「実務限に関する様々な立場について」\cite{cantor}および\cite{kmoonshine}1.8節を参照されたい，

        \subsection{p.170, l.16-}

            \begin{quote}
                「あまりにも素朴な万民動員型の闘争の結果であった〜」
            \end{quote}

            ホッブズの「万人の万人に対する闘争」\cite{hobbes1}に由来する文言．\ref{malthus}節で見た人口爆発を受けて，人口削減のために大規模な殺戮が発生したということ．ホッブズは『物体論』\cite{hobbes2}において，数学なしには正しい自然認識はありえないと主張していることに注意．

            また，村上陽一郎『新しい科学論』\cite{murakami}p.193に「万人の万人に対する闘い」という文言が現れている．これは原子論が西洋社会における近代人としての自我の確立に対応する，つまり，社会における最小構成単位である個人の確立と，物質の最小構成単位である（と当時は考えられていた）原子論の確立は対応する，という文脈で登場する．

        \subsection{p.171, l.7}

            \begin{quote}
                「プロジェクト・ヘロデ」
            \end{quote}

            新約聖書に由来する．新約聖書において，自らの支配地に“ユダヤ人の王”（イエス・キリストのこと）が生まれると聞いたユダヤ王ヘロデは，自らの支配地に生まれる男児を皆殺しにすることを決意し，実行した．イエスの母マリアの夫であるヨセフは，予言に従ってマリアと共にエジプトに逃れ，ことなきを得た\cite{bible}．

        \subsection{p.172, l.2-}

            \begin{quote}
                「こうしてあなたは六つ目の断章へと辿りつき，〜」
            \end{quote}

            文中の指示に従ってきちんと読んできただろうか．

        \subsection{p.172, l.12-p.173, l.2}

            \begin{quote}
                「あなたは私を振り返り，〜」
            \end{quote}

            2004の語り手は“私”．

        \subsection{p.173, l.9-}

            \begin{quote}
                「要するに黒の数から白の数を差っ引いて，〜」
            \end{quote}

            チューリング機械の動作を直接に描写している．テクストを読んで状態を記録し，前後のテクストに移動するアルゴリズムとは，チューリング機械にほかならない．直後に有限チューリング機械の話題が登場することから，この見立ては十二分に支持される．

            また，白と黒が連続すること，その差分を取ることから，東大の碁石の問題として知られる，有名な入試問題が連想される．

            \begin{prb}{東大文系第4問(2001)}{}
                白石180個と黒石181個の合わせて361個の碁石が横に一列に並んでいる．碁石がどのように並んでいても，次の条件を満たす黒の碁石が少なくとも1つあることを示せ．
                \begin{itemize}
                    \item その黒の碁石とそれより右にある碁石をすべて除くと，残りは白石と黒石が同数となる．
                \end{itemize}
                ただし，碁石が1つも残らない場合も同数とみなす．
            \end{prb}

            様々な解法が考えられるが，標準的な解法は中間値の定理のような論法に持ち込むものであろう\footnote{本書は受験参考書ではないので解法は省略する．}．なお，本作の状況は東大の碁石の問題とは異なることに注意せよ．この問題は2000年代後半以降の高校生には有名かもしれないが，円城塔が知らない可能性は十分に考えられる．

        \subsection{p.174, l.15}

            \begin{quote}
                「専門家の間では言語のクラスと呼ばれている」
            \end{quote}

            前後の文脈から，計算機科学における言語を指していることがわかる．ここで，計算機科学における言語は以下のように定義される\cite{sipser1}．

            \begin{dfn}{（計算機科学における）言語，言語のクラス}{}
                機械$M$が受理するすべての文字列の集合$A$を，機械$M$の\textbf{言語}であるという．また，言語の集まりのことを\textbf{クラス}という．
            \end{dfn}

            言語は集合であるから，“言語の集合”を素朴に定義することは出来ず，言語のクラスという形で定義されている．作中の描写より，“その種の限定”はチューリング機械のテープが有限であることを指しているから，当該箇所は，どのオートマトンに受理されるかによって言語のクラスが分類されることを指している．

        \subsection{p.176, l.8-10}

            \begin{quote}
                「実際にはもう少し面倒くさい細部もあるわけであり，〜」
            \end{quote}

            ここに挙げられている問題意識は，のちに『文字渦』で集中して扱われる．

        \subsection{p.176, l.11}

            \begin{quote}
                「あなたの好きに読まれている．」
            \end{quote}

            自らの作品が，自らの意図した通りに読まれないことをはじめから割り切っているとも読める．『文字渦』以降は意図通りに読まれないことを前提としている感が強いが，最初期はまだある程度意図通り読んでもらえないものかと試行錯誤している気配がある．

        \subsection{p.178, l.1-2}

            \begin{quote}
                「ひとつの事故とひとつの事件と，〜」
            \end{quote}

            2005の語り手は無人称．

        \subsection{p.179. l.6-7}

            \begin{quote}
                「人命はガスジャイアント程にも浮薄であり，」
            \end{quote}

            ガスジャイアントとは巨大ガス惑星のこと．太陽系では木星と土星が該当する．いずれも比重が非常に軽いことで有名であり，特に土星は比重が1より小さい，つまり同じ体積の水よりも軽いことが知られている．これは“人命は地球より重い”という名言を前提とした皮肉である．元々はサミュエル・スマイルズ『自助論』（『西国立志編』）\cite{smiles}に見え，わが国においては1948年の死刑合憲判決\cite{justice}や1977年のダッカ日航機ハイジャック事件に対する福田赳夫首相の発言に引用されたことで広く知られている．

            科学者を狙ったテロはグレッグ・イーガン『万物理論』からか．

        \subsection{p.179, l.10-11}

            \begin{quote}
                「奴らが人類のために生み出した最善のものはと訊ねられて，原子爆弾という答えが返る〜」
            \end{quote}

            ジョン・フォン・ノイマンは強硬なタカ派として知られており，特に「もし奴らを明日爆撃するというならば，私はなぜ今日でないのかと言う．今日の5時に爆撃するというならば，私はなぜ1時ではないのかと言う．」と語ったとされること\cite{life}で有名である．フォン・ノイマンは第二次大戦中はマンハッタン計画に参加していたことから，このような危険思想を持つ科学者が実際に原爆開発に携わっていた事実を元にした記述か．

            あるいは，村上陽一郎『科学者とは何か』\cite{murakami}において，ファインマン『ご冗談でしょう，ファインマンさん』\cite{fynman1}が悪意をもって引用され，マンハッタン計画に参加していた物理学者が軒並み不当に気狂い扱いされていたこと\footnote{詳細については後日Appendixとして述べる．}を揶揄する記述か．

        \subsection{p.181, l.1-2}

            \begin{quote}
                「片手を上げて厳かに静粛を求めた次の瞬間，この地域一体が巨大な地震に襲われたところにある．」
            \end{quote}

            2004年12月に発生したスマトラ島沖地震をもとにした描写か．

        \subsection{p.182, l.2-3}

            \begin{quote}
                「二〇〇五年三月，同地域を偶然にも連続して壊滅させることになる大規模な地震」
            \end{quote}

            2005年3月に発生したスマトラ島沖地震のこと．有名なのは前者の2004年発生のもの．

        \subsection{p.182, l.14}

            \begin{quote}
                「会議名は，普遍計算と力学系と銘打たれており，〜」
            \end{quote}

            実際にこのような会議が開催されたかは不明．「Boy's Surface」に書かれている通り，科学者の国際会議に関する記述は基本的に信用ならない
                %\ref{meeting}節を参照のこと．
            ．

            普遍計算はuniversal conputationのこと．あらゆることは計算であり，あらゆることは力学系であるから，この会議はあらゆることを対象としたものであるというまったく情報量のない洞察が得られる．

        \subsection{p.183, l.11}

            \begin{quote}
                「ユニバース．」
            \end{quote}

            無論，宇宙のことだが，複数の意味を同時に象徴しているか．

            まず，直後に“集合の集合”についての言及があることから，集合論や数理論理学で用いられる宇宙が考えられる．ここで，集合論における宇宙とは，以下のように定義される\cite{arai,tanakak}．

            \begin{dfn}{（集合論における）言語，構造，宇宙}{}
                関数記号（$+, \,\cdot\,, 0, 1, \cdots$），関係記号（$=, >, \in, \cdots$）から成る集合Lを\textbf{言語}(language)という．

                $M$を空でない集合，$F$をLの各記号$\alpha$について以下のように$\alpha^\mathcal{M} = F(\alpha)$を対応させる写像とする．

                \begin{itemize}
                    \item $m$-変数関係記号$R \in \textrm{L}$に対して$R^\mathcal{M}$は$M$上の$m$-項関係，すなわち$R^\mathcal{M} \subset M^m$．
                    \item $n$-変数関数記号$f \in \textrm{L}$に対して$f^\mathcal{M}$は$M$上の$n$-項関数，すなわち$f^\mathcal{M} : M^n \to M$．
                    \item 特に定数$c \in \textrm{L}$に対して$c^\mathcal{M} \in M$．
                \end{itemize}

                このとき，言語Lにおける\textbf{構造}(structure)とは，対$\mathcal{M} = \langle M ; F \rangle$である．$\mathcal{M}$を通常$\langle M ; R^\mathcal{M}, \cdots, f^\mathcal{M}, \cdots, c^\mathcal{M}, \cdots \rangle$と書き，$M = |\mathcal{M}|$を構造$\mathcal{M}$の\textbf{宇宙}(universe)，あるいは\textbf{対象領域}，または単に\textbf{領域}という．
            \end{dfn}

            集合論における宇宙とは異なる概念で，集合論の周辺で用いられる宇宙も存在する．

            \begin{dfn}{遺伝的有限集合}{}
                \textbf{遺伝的有限集合}とは，以下の条件をすべて満たす有限集合である．

                \begin{itemize}
                    \item 空集合$\varnothing$は遺伝的有限集合である．
                    \item 遺伝的有限集合の冪集合は遺伝的有限集合である．
                \end{itemize}
            \end{dfn}

            \begin{dfn}{グロタンディーク宇宙}{}
                \textbf{グロタンディーク宇宙}とは，以下の条件をすべて満たす集まり$U$である．

                \begin{itemize}
                    \item $x \in U, y \in x$ならば，$y \in U$
                    \item $x, y \in U$ならば，$\{ x, y \} \in U$
                    \item $x \in U$ならば，$P(x) \in U$
                    \item $\{ x_\alpha \}_{\alpha \in I}$が$U$の元の族で$I \in U$ならば，$\bigcup_{\alpha \in I} x_\alpha \in U$
                \end{itemize}
            \end{dfn}

            グロタンディーク宇宙とは，平たく言えば，通常の集合論で使い得る数学的対象全てから成る“集合”である\footnote{グロタンディーク宇宙は“集合の集合”であるため，公理集合論においては集合ではない．このことは，「ユニバース．」がグロタンディーク宇宙をも想定していることを強力に示唆する．}．それさえ用意しておけば，その内部で通常の数学をすべて展開可能であるという，ある種理想的な環境のことをグロタンディーク宇宙と呼んでいる．グロタンディーク宇宙には階層が存在し，最小のグロタンディーク宇宙は空集合$\varnothing$である\footnote{$\varnothing$は自明なグロタンディーク宇宙であるともに，自明なクワインでもある．これを直接の題材にしたのが「$\varnothing$」．}．次に小さなグロタンディーク宇宙はすべての遺伝的有限集合$V_\omega$であり，この中で通常の数学はすべて表現可能．

            次に，様相論理におけるクリプキ意味論の宇宙が挙げられる．ここで，様相論理におけるクリプキ意味論およびそれを適用すべきクリプキフレームとクリプキモデルは以下のように定義される\cite{yoso}．

            \begin{dfn}{クリプキフレーム，クリプキモデル}{}
                \textbf{クリプキフレーム}とは，空でない集合$W$と$W$上の2項関係$R \subseteq W \times W$の対$(W, R)$である．なお，クリプキ意味論では，$W$の要素$w$は世界（宇宙）に相当し，$R$は世界間の遷移関係・到達可能性に相当する．

                また，\textbf{クリプキモデル}とは，クリプキフレーム$(W, R)$と，$W$上の付値関数$V: \textsf{Prop} \to \mathcal{P}(W)$の組$(W, R, V)$である．
            \end{dfn}

            また，宇宙論における宇宙，つまり通常の意味での宇宙でもある．

            さらに，直後に登場する万能チューリング機械(Universal Turing Machine)の示唆でもある．

            なお，ABC予想の解決に用いられるとされる宇宙際タイヒミュラー理論における宇宙はこれらとはまた異なる概念であり，しかも本作の描写には不適．

        \subsection{p.183, l.12}

            \begin{quote}
                「可能なものの全ての集合は集合を超えてしまっている．」
            \end{quote}

            正しい．“可能なものの全ての集合”を集合であると認めてしまうと，ラッセルの逆理が成り立ってしまう．これを回避するため，公理的集合論では“集合の集合”は集合族として定義されている．詳細は\cite{kmoonshine}Appendix Aを見よ．

        \subsection{p.184, l.1-4}

            \begin{quote}
                「チューリング当人による構成におけるユニバーサル・チューリング・マシンは，〜」
            \end{quote}

            万能チューリング機械に関する説明．簡単のため正確性は犠牲になっているが，間違ったことは言っていない．

        \subsection{p.184, l.5-11}

            \begin{quote}
                「チューリングマシンに関する簡潔な説明をお望みだろうか．〜」
            \end{quote}

            前半はチューリング機械，後半はチャーチ\hyphen チューリングの提唱に関する簡潔な説明．いずれも正しい．

        \subsection{p.184, l.12-16}

            \begin{quote}
                「私は，人の怠惰と，〜」
            \end{quote}

            テープ\hyphen ヘッドから成る構成，つまりチューリング機械によって現象を解明しようという姿勢は，多くの作品に見られる姿勢である．円城塔の科学的思考における常套手段とも言える．

            本作に関しては，大元はチューリング機械だが，既に述べている通り池上高志の池上機械からの影響も指摘される．池上機械の詳細については\cite{tanakah}を参照されたい．

        \subsection{p.185, l.13-}

            \begin{quote}
                「私を読み出したあなたの中に生ずる印象は，私よりも真に大きなものでありえて，〜」
            \end{quote}

            正しい．ダグラス・アダムス『銀河ヒッチハイク・ガイド』における，地球は実は計算機であったのだが，それを地球上の生命は誰も知らなかった，というギャグにも通じる．

        \subsection{p.186, l.3-5}

            \begin{quote}
                「自己認識ニューロンの活動が，運動を企図するニューロンの準備電位に遅れていることが知られたのは，恐竜がまだ地表を闊歩していた一九七〇年代の話である．」
            \end{quote}

            ベンジャミン・リベットによる有名な実験のことを指すと推定されるが，リベットの実験は1980年代に実施されたものであり，齟齬が生じる．後半の「恐竜がまだ〜」はギャグ．

        \subsection{p.186, l.12}

            \begin{quote}
                「禁止集合」
            \end{quote}

            ある問題について考えるとき，可能ではあるが除外して考えたい，つまり禁止したい状態の集まりを禁止集合ということがある．組合せ論，グラフ理論をはじめ，数理科学全般で使いうる表現である．ただし，私の記憶では，物理学で用いている例を目にしたことはない．

            この箇所の文学的な解釈は定まらない．禁止集合を利用する曖昧な計算とはどのようなものだろうか．検知不能な不可知領域を間借りして，自分が存在するための計算を行なっている，ということか？

        \subsection{p.187, l.1}

            \begin{quote}
                「それでもあなたに，私を自由に読むという権利が与えられていることは自明であり，〜」
            \end{quote}

            小説というメディアは，読み手に最大限の自由度を与える．私はいま，本作の註釈を与えているが，読者はこれを受け入れてもいいし，当然すべて拒否してもよい．

            円城塔は，読者に対して，読者それぞれで読み方は自由で開かれているとするが，初期作品においては，自分の想定したように読まれなかったことへの驚きと微かな失望を見せることがある．

            もっとも極端な例が，（初期と中期の境界ではあるが）『プロローグ』にある，「これはペンです」が芥川賞の選考の際に村上龍から科学的にまったく誤った知識に基づく誹謗中傷を受けたことに対する恨みの文言\footnote{「こういうことを書いていると，それは生物学の知識として間違っている，そういうところが駄目だと小説も駄目になるんですよ，僕は昔そのあたりのことを調べたことがあるからよく知っているんですとか，したり顔で言い出す人が出てくるが，まあ無視してよろしい．」\cite{prologue}p.298}である．これについては，円城塔にまったく同情する．

            他にも，円城塔作品に対して寄せられる心無い誤解を私が多く目にしたこと，円城塔からの働きかけに答える読者がまったく見られないこと，円城塔作品に対してあまりにも酷い無理解が蔓延していることなどが，私がこの註釈を書く動機となった．

            少なくとも，作品を通じて読者に対して行なったはずの働きかけに対して，読者からまったく応答がないと，結構露骨にガッカリしているので，そこに対してきちんと応答していきたい所存．

        \subsection{p.189, l.6}

            \begin{quote}
                「禍福はアズラエルの罠の如し」
            \end{quote}

            「禍福はあざなえる縄のごとし」のもじり．

        \subsection{p.189, l.14}

            \begin{quote}
                「ゼブラガニ」
            \end{quote}

            実在の蟹．ラッパウニ，イイジマフクロウニに共生し，棘や管足を切り取って食べるほか，そのウニを操縦して移動までこなす．したがって，作中の描写は正しい．名前は外見が白と濃褐色の縦縞模様であることから．

        \subsection{p.190, l.12-13}

            \begin{quote}
                「体全体の情報を統合させずとも，大過なく暮らしていける構成を持っている」
            \end{quote}

            ウニ，ヒトデについての記述．これは正しい．ウニ，ヒトデなどの棘皮動物は，脳のような明確な中心処理系が見られず，体内の神経系による分散制御を行なっていると考えられている\cite{takahashi}．

        \subsection{p.195, l.9-10}

            \begin{quote}
                「ウイルスに対する完璧なセキュリティソフトが存在しないことの証明にもまた，ユニバーサルチューリングマシンが用いられる〜」
            \end{quote}

            正しい．コーエン\footnote{このフレッド・コーエンは，強制法を開発した数学者ポール・コーエンとは別人．余談だが，コンピュータウイルスという語を学術的に初めて用いたのは，フレッド・コーエンの指導教官であるレナード・エイドルマン（RSA暗号のA）とのこと\cite{cohen}．}による証明\cite{cohen}がある．

        \subsection{p.195, l.14-15}

            \begin{quote}
                「あなたを十重二十重に包囲しているマシンたち〜」
            \end{quote}

            ここから，スサノオの短歌「八雲立つ出雲八重垣妻籠みに八重垣つくるその八重垣を」\cite{kojiki}が連想される．この短歌は「Goldberg Invariant」で明示的に登場する%（\ref{susanoo}節を参照のこと．）
            ことから，飛躍ではないと考えられる．

        \subsection{p.195, l.17-}

            \begin{quote}
                「鏡と生殖は存在の数を無闇と増やすために邪悪であり，〜」
            \end{quote}

            ホルヘ・ルイス・ボルヘス「トレーン，ウクバール，オルビス・テルティウス」\footnote{「鏡と交合は人間の数を増殖するがゆえにいまわしい」\cite{fictions}p.14}をほんの少し改変したもの．

        \subsection{p,197, l.1-4}

            \begin{quote}
                「とここに記しておいて，〜」
            \end{quote}

            テクストの記述に従って，地球の半径を6400\,kmとおく．実際に計算してみると，1人あたりの面積は約73500\,$\mathrm{m}^2$となり，確かに1辺300\,mの正方形の面積90000\,$\mathrm{m}^2$を割り込んでいることがわかる．

        \subsection{p.198, l.13-14}

            \begin{quote}
                「実際僕は，こうして苦肉の策としてあいつの企みに乗っかり利用しようとする代わりに，何がしかのものを差し出してしまっているわけであり，〜」
            \end{quote}

            本作には，“僕”“私”そして“無人称”の3人の語り手がいる．ここで“あいつ”と呼ばれている存在が“無人称”である．

        \subsection{p.199, l.1-2}

            \begin{quote}
                「〜どことも知れぬ浜辺を目指す．そんな馬鹿げた試みを心底愛し，今も継続し続けている決して零ではない数の人々」
            \end{quote}

            科学者，特に複雑系の研究者のことを指すと推察される．物理学者の田崎晴明は，90年代後半の複雑系の研究者の一部に見られたという，計算機や自然言語，安易な思いつきに過度に依存した態度に対して，次のように批判した\cite{tazaki}．

            \begin{quote}
                本当の「複雑系」への道はものすごく険しいもの． 結局，何一つ普遍的なもののみつからない砂漠をさまよい続けて，討ち死にする覚悟がなくては旅立てないような世界のはずだ． 計算器がまわせれば参入できるお手軽な流行の分野というとらえ方があるのは，全くおかしい．
            \end{quote}

            あてもなく浜辺を彷徨う人々，というモチーフは，この複雑系への批判に由来しているように思われた．本作が複雑系の研究者時代の自伝に近い作品であることからも，この見立ては十分に支持される．

	\begin{thebibliography}{99}

		\bibitem{boys} 円城塔, 『Boy's Surface』, ハヤカワ文庫JA, 早川書房, 2011

        \bibitem{fujita} 藤田直哉, 科学と文学は結婚できるのか？ : 円城塔「Your Head's Only」論, \textit{思辨改}, 2009, \url{http://speculat.cocolog-nifty.com/blog/2009/02/your-heads-only.html}

        \bibitem{chaos} 金子邦彦, 『カオスの紡ぐ夢の中で』, ハヤカワ文庫NF, 早川書房, 2010

        \bibitem{sipser2} Michael Sipser, 『計算理論の基礎2』, 原書第3版, 共立出版, 2023

        \bibitem{mumonkan} 無門慧開, 『無門関』, 岩波文庫, 岩波書店, 1994

        \bibitem{kimura} 木村達雄, 数学は体力だ！, \textit{筑波フォーラム}, (45), 104-107, 1996

        \bibitem{sri} 円城塔, 『Self-Referential INTERVIEW』, 2007, \url{https://engine.sub.jp/ppt/SRI.ver2.3.pdf}

        \bibitem{bunshun1} 宮本輝, さざ波と小技, \textit{文藝春秋}, \textbf{89}(11), 2011

        \bibitem{bunshun2} 島田雅彦, 人を虚仮にしたっていいじゃないか，小説だもの, \textit{文藝春秋}, \textbf{90}(4), 2012

        \bibitem{nks} Stephen Wolfram, ``A New Kind of Science'', 2002, \url{https://www.wolframscience.com/nks/}

        \bibitem{cook} Matthew Cook, Universality in elementary celluar automata, \textit{Complex Systems}, \textbf{15}(1), 1-40, 2004, \url{https://doi.org/10.25088/ComplexSystems.15.1.1}

        \bibitem{wing} Lorna Wing, Asperger's syndrome : a clinical account, \textit{Psychological Medicine}, \textbf{11}(1), 115-129, 1981, \url{https://doi.org/10.1017/S0033291700053332}

        \bibitem{hongo} 本郷一夫, 八木成和, 『シードブック教育心理学』, 建帛社, 2008

        \bibitem{buchen} Lizzie Buchen, Scientists and autism : when geeks meet, \textit{Nature}, \textbf{479}, 25-27, 2011, \url{https://doi.org/10.1038/479025a}

        \bibitem{dirac} アブラハム・パイス, 『ポール・ディラック』, ちくま学芸文庫, 筑摩書房, 2012

        \bibitem{interview1} 円城塔, 円城塔さんインタビュー, \textit{Book Shorts}, 2015, \url{https://bookshorts.jp/enjoetoh/}

        \bibitem{interview2} 円城塔, 佐々木敦, 『エピローグ』と『プロローグ』のあいだ : 世界・SF・私小説, \textit{ECRIT-O}, 2018, \url{https://ecrito.fever.jp/20180701223008}

        \bibitem{savarese} ラルフ・ジェームズ・サリヴァーズ, 『嗅ぐ文学，動く言葉，感じる読書』, みすず書房, 2021

        \bibitem{sasaki} 佐々木敦, 『あなたは今，この文章を読んでいる．』, 慶應義塾大学出版会, 2014

        \bibitem{posdoc} 円城塔, ポスドクからポストポスドクへ, \textit{日本物理学会誌}, \textbf{63}(7), 564-566, 2008, \url{https://doi.org/10.11316/butsuri.63.7_564}

        \bibitem{kanamori} 金森修, 『サイエンス・ウォーズ』, 新装版, 東京大学出版会, 2014

        \bibitem{objectibity} ロレイン・ダストン, ピーター・ギャリソン, 『客観性』, 名古屋大学出版会, 2021

        \bibitem{fabre} ファーブル, 『完訳ファーブル昆虫記2』, 岩波文庫, 岩波書店, 1993

        \bibitem{nakane} 中根美知代, 『$\varepsilon$--$\delta$論法とその形成』, 共立出版, 2010

        \bibitem{boyers} David G. Boyers, William A. Tiller, Corona discharge photography, \textit{Journal of Applied Physics}, \textbf{44}(7), 1973, \url{https://doi.org/10.1063/1.1662715}

        \bibitem{fukurai} 福来友吉, 『透視と念写』, 東京宝文館, 1913

        \bibitem{fuji} 藤教篤, 藤原咲平, 『千里眼実験録』, 大日本図書, 1911

        \bibitem{emoto} 江本勝, 『水からの伝言』, 波動教育社, 1999

        \bibitem{amabane} 甘羽優子，田崎晴明，菊池誠，「水からの伝言」をめぐって，\textit{日本物理学会誌}, \textbf{66}(5), 342-346, 2011, \url{https://doi.org/10.11316/butsuri.66.5_342}

        \bibitem{cell} 金子邦彦, 澤井哲, 高木拓明, 古澤力, 『細胞の理論生物学』, 東京大学出版会, 2020

        \bibitem{lolita} ウラジーミル・ナボコフ, 『ロリータ』, 新潮文庫, 新潮社, 2006

        \bibitem{waki} 脇司, 尾針由真, 林慶, 森部絢嗣, 松尾加代子, 高島康弘, 反芻動物を宿主とする槍形吸虫類の一種Dicrocoelium chinensisの日本におけるスポロシストの初検出, \textit{日本野生動物医学会誌}, \textbf{27}(1), 35-43, 2022, \url{https://doi.org/10.5686/jjzwm.27.35}

        \bibitem{cantor} Georg Cantor, ``Gesammelte Abhandlungen mathematischen und philosophischen Inhalts'', 1932, Springer, \url{http://resolver.sub.uni-goettingen.de/purl?PPN237853094}

        \bibitem{kmoonshine} 下村思游, 『円城塔全註釈12 ムーンシャイン』, 2025, 柏屋

        \bibitem{hobbes1} ホッブズ, 『リヴァイアサン2』, 改訳版, 岩波文庫, 岩波書店, 1992

        \bibitem{hobbes2} ホッブズ, 『物体論』, 京都大学学術出版会, 2015

        \bibitem{murakami} 村上陽一郎, 『新しい科学論』, 講談社ブルーバックス, 講談社, 1979

        \bibitem{bible} 『文語訳新約聖書』, 岩波文庫, 岩波書店, 2014

        \bibitem{sipser1} Michael Sipser, 『計算理論の基礎1』, 原書第3版, 共立出版, 2023

        \bibitem{smiles} サミュエル・スマイルズ, 『西国立志編』, 講談社学術文庫, 講談社, 1981

        \bibitem{justice} 最高裁判所大法廷, 昭和22(れ)119, \textit{最高裁判所刑事判例集}, \textbf{2}(3), 191, 1948, \url{https://www.courts.go.jp/assets/hanrei/hanrei-pdf-56385.pdf}

        \bibitem{life} Clay Blair Jr., Passing of a Great Mind, \textit{Life}, \textbf{42}(8), 89-104, 1957

        \bibitem{murakami} 村上陽一郎, 『科学者とは何か』, 新潮選書, 新潮社, 1994

        \bibitem{fynman1} リチャード・P・ファインマン, 『ご冗談でしょう，ファインマンさん 上』, 岩波現代文庫, 岩波書店, 2000

        \bibitem{arai} 新井敏康, 『数学基礎論』, 増補版, 東京大学出版会, 2021

        \bibitem{tanakak} 田中一之, 『数学基礎論序説』, 裳華房, 2019

        \bibitem{yoso} 佐野勝彦, 倉橋太志, 薄葉季路, 黒川英徳, 菊池誠, 『数学における証明と論理 : 様相論理と数学基礎論』, 共立出版, 2016

        \bibitem{tanakah} 田中博, 『生命と複雑系』, 培風館, 2002

        \bibitem{prologue} 円城塔, 『プロローグ』, 文春文庫, 文藝春秋, 2018

        \bibitem{takahashi} 髙橋景一, ウニの棘の運動, \textit{比較生理生化学}, \textbf{14}(3), 185-193, 1997, \url{https://doi.org/10.3330/hikakuseiriseika.14.185}

        \bibitem{cohen} Fred Cohen, Computer viruses : theory and experiments, \textit{Computers \& Security}, \textbf{6}(1), 22-35, 1987, \url{https://doi.org/10.1016/0167-4048(87)90122-2}

        \bibitem{kojiki} 『古事記』, 改版, 岩波文庫, 岩波書店, 2007

        \bibitem{fictions} ホルヘ・ルイス・ボルヘス, 『伝奇集』, 岩波文庫, 岩波書店, 1993

        \bibitem{tazaki} 田崎晴明, 「黒木のなんでも掲示板」における田崎の掲示 : 「複雑系」と計算器, 1999, \url{https://haltasaki.github.io/writings/html/kuroki/comp.html}

    \end{thebibliography}

\end{document}